%% =====================================================================
%%  T-07-02  The Indirect Question
%%  -------------------------------------------------------------------
%%  One idea per file. Core is mandatory. Slots in canonical order.
%%  Max 700 words. No layout commands. No filler. New source material
%%  for this entry goes in sources/<BOOK>/T-07-02.tex -- never by
%%  rewriting this file (docs/RULES.md R0.1).
%% =====================================================================
%% @kind: topic
%% @id: T-07-02
%% @title: The Indirect Question
%% @subtitle: Ask beside the thing you want to know
%% @chapter: c07
%% @order: 20
%% @status: stable
%% @sources: B3, B1
%% @updated: 2026-09-19

\topic{T-07-02}{The Indirect Question}{Ask beside the thing you want to know}

\begin{core}
A direct question about a sensitive subject produces a prepared answer. The same information obtained from an adjacent question --- about a similar case, a third party, a hypothetical --- produces an unprepared one.
\end{core}
\begin{mechanism}
Preparation is topic-specific. A person who has thought about how to answer whether they intend to leave has not thought about how to answer what they would do if someone in their position left. The adjacent question routes around the prepared answer while keeping the underlying disposition engaged, and it does so without triggering the defensiveness that a direct approach causes. Two variants work particularly well: asking about a third party, which lets the respondent project their own view safely, and asking for advice, which invites them to state the reasoning they would otherwise conceal.
\end{mechanism}
\begin{tells}
  \item Someone asks your opinion about a case that resembles your situation closely.
  \item You are asked what you would do, or what you advise, about something that concerns them.
  \item A question about a third party has an oddly specific shape.
  \item The asker reacts to your answer rather than continuing the conversation.
  \item Questions cluster around one subject across several unrelated conversations.
\end{tells}
\begin{moves}
  \item Move one step sideways: ask about someone like them, someone near them, or a hypothetical.
  \item Ask for advice on a decision that is theirs. Advice-givers reveal their own criteria.
  \item Ask about a third party's comparable situation and let them generalise.
  \item Follow the reaction, not the words: hesitation and correction are more informative than content.
  \item Stop after one indirect question on a subject. Two is a pattern they will notice.
\end{moves}
\begin{counters}
  \item Notice the sideways step. When a question is adjacent to something sensitive, answer the adjacent question and not the underlying one.
  \item Answer hypothetically in kind: give a general view that commits you to nothing specific.
  \item Reverse it --- ask what prompted the question. Genuine curiosity survives that question; collection does not.
  \item Track clustering. A single indirect question is conversation; several around one subject across time is a campaign.
\end{counters}
\begin{cost}
Indirection produces indirect information. The respondent may answer the adjacent question honestly while the real position stays hidden, and projecting a view onto a third party is not the same as holding it. There is also a reciprocity cost: people who are habitually asked sideways eventually notice, and the noticing damages the relationship the information came from. The method is strongest on dispositions and weakest on facts, which are better obtained from documents.
\end{cost}

\sources{B3, B1}
\seesources
\seealso{T-07-01, T-07-03, T-05-01}
